% ============================================================
% Introduction
% ============================================================
\subsection{Cartilage, the growth plate, and the extracellular matrix}
\label{subsec:cartilage-growth-plate-ecm}

Longitudinal bone growth in vertebrates occurs through endochondral ossification, a process in which a cartilage template is progressively expanded, remodeled, and replaced by bone. The growth plate is therefore not only a cellular differentiation system, but also an extracellular matrix-rich tissue whose mechanical and signalling properties depend on the composition of its matrix.
Cartilage extracellular matrix is dominated by fibrillar collagens, especially collagen type II, and by proteoglycans that contribute to hydration, tissue resilience, and the spatial organization of matrix components.

\subsection{Growth-plate zonation and a core non-SLRP gene network}
\label{subsec:growth-plate-nonslrp-network}

The epiphyseal growth plate is organized into resting, proliferative,
prehypertrophic, and hypertrophic chondrocyte compartments. These zones are
not separate tissues, but successive cellular states whose coordinated
proliferation, differentiation, matrix production, and matrix removal produce
longitudinal growth. Several non-SLRP genes provide the regulatory and matrix
framework within which the SLRPs investigated in this thesis act
(Table~\ref{tab:growth-plate-nonslrp}).

\begin{table}[H]
\centering
\caption{Representative non-SLRP genes in growth-plate organization. The
selection summarizes major regulatory and matrix modules and is not intended
as an exhaustive catalogue of all genes expressed in growth-plate cartilage.}
\label{tab:growth-plate-nonslrp}
\begin{tabularx}{\textwidth}{@{}p{0.27\textwidth}X@{}}
\toprule
Gene or pathway & Principal growth-plate role \\
\midrule
\textit{SOX9}, \textit{COL2A1}, and \textit{ACAN} &
SOX9 is essential for chondrocyte lineage formation and supports expression of
cartilage-matrix genes including type-II collagen and aggrecan. These matrix
components provide the collagenous and hydrated proteoglycan framework in
which SLRPs are embedded \citep{Bi1999Sox9}. \\
\textit{IHH}, \textit{PTHLH}, and \textit{PTH1R} &
The IHH--PTHrP feedback system couples the position of prehypertrophic cells to
the rate at which proliferating chondrocytes enter hypertrophy, thereby
preventing premature terminal differentiation \citep{Vortkamp1996IhhPthrp}. \\
\textit{RUNX2} and \textit{COL10A1} &
RUNX2 is required for normal chondrocyte proliferation and maturation during
endochondral ossification. Type-X collagen is a characteristic component of
hypertrophic cartilage and marks the matrix state that precedes replacement by
bone \citep{Chen2014Runx2}. \\
\textit{FGFR3} &
FGFR3 limits endochondral bone growth; disruption of the receptor in mice
expands proliferating and hypertrophic growth-plate compartments, whereas
activating human variants cause skeletal dysplasias such as achondroplasia
\citep{Deng1996Fgfr3}. \\
\textit{NPPC} and \textit{NPR2} &
Locally produced C-type natriuretic peptide signals through NPR2/guanylyl
cyclase-B and positively supports endochondral ossification and longitudinal
bone growth \citep{Tamura2004Npr2}. \\
\textit{MMP13} &
MMP13 cleaves cartilage matrix during the hypertrophic-to-ossification
transition. Loss of MMP13 causes collagen accumulation, enlargement of the
hypertrophic zone, and delayed ossification and vascularization
\citep{Inada2004Mmp13}. \\
\textit{VEGFA} &
VEGF produced by hypertrophic chondrocytes promotes vascular invasion and
couples terminal cartilage remodeling to trabecular bone formation
\citep{Gerber1999Vegf}. \\
\bottomrule
\end{tabularx}
\end{table}

This framework separates different kinds of evidence. SOX9, IHH--PTHrP,
RUNX2, FGFR3, and CNP--NPR2 primarily regulate chondrocyte state or growth;
COL2A1, ACAN, and COL10A1 define changing matrix environments; and MMP13 and
VEGFA help execute the cartilage-to-bone transition. The focal SLRPs are best
placed within the matrix branch of this system: their conserved secretion
signals, leucine-rich repeat architecture, and collagen-associated functions
make them plausible modifiers of the extracellular environment, but do not by
themselves establish that an SLRP controls growth-plate zonation. These
non-SLRP genes therefore provide biological context rather than an additional
comparative candidate set in the present study.

\subsection{Proteoglycans and small leucine-rich proteoglycans}
\label{subsec:proteoglycans-slrps}

Small leucine-rich proteoglycans (SLRPs) are extracellular matrix proteins characterized by a relatively small core protein and a central leucine-rich repeat domain. This domain forms a curved interaction surface that enables binding to collagens, growth factors, and other matrix-associated molecules. Although historically described mainly as structural matrix components, SLRPs are now understood as multifunctional regulators of matrix assembly, tissue mechanics, mineralization, and cell signalling \citep{SchaeferIozzo2008,IozzoSchaefer2015}.

\subsection{SLRPs in collagen-rich skeletal tissues}
\label{subsec:slrps-skeletal-tissues}

SLRPs influence collagen-rich tissues by binding to fibrillar collagens and modulating fibril assembly, spacing, and susceptibility to proteolytic degradation. This is particularly relevant in cartilage, where matrix remodeling accompanies both normal skeletal development and pathological degeneration. For example, fibromodulin is processed by cartilage-associated proteases such as MMP-13 and ADAMTS enzymes, and MMP-13-cleaved fibromodulin has been detected in hypertrophic chondrocytes of human fetal growth plates.

\subsection{Evolutionary origin and species sampling}
\label{subsec:slrp-evolutionary-origin}

Comparative studies indicate that an ancestral SLRP-like protein was probably
present early in chordate evolution and that the family expanded during early
vertebrate evolution. Tunicates such as \textit{Ciona} contain a small number
of SLRP-like genes that have been discussed as possible precursors of the
canonical vertebrate family, whereas early-diverging vertebrates already show
several distinguishable SLRP-like genes \citep{Park2008}. A broadly recognizable
modern repertoire was established before the divergence of cartilaginous and
bony vertebrates, followed by lineage-specific changes including teleost gene
duplication \citep{Costa2018}. A more recent comparative survey describes
both whole-genome and tandem duplication as contributors to the expansion of
clustered SLRP genes in jawed vertebrates \citep{Gil2025}.

No living species is treated here as the ancestor from which the other sampled
species evolved. Instead, amphioxus, lamprey, cartilaginous fishes, spotted gar,
coelacanth, teleosts, and tetrapods bracket informative stages of chordate and
vertebrate history. \textit{Ciona} is retained as literature-only family-origin
context because its SLRP-like sequence cannot be assigned confidently to BGN,
DCN, FMOD, PRELP, EPYC, LUM, or OGN. Amphioxus defines a deep chordate
uncertainty boundary; lamprey tests early vertebrate diversification; sharks
represent the earliest sampled jawed-vertebrate branch; and gar and coelacanth
help distinguish ancient vertebrate conservation from teleost-specific
duplication.

\subsection{Candidate genes investigated in this thesis}
\label{subsec:candidate-genes}

The comparative panel comprises \textit{BGN}, \textit{DCN}, \textit{FMOD},
\textit{PRELP}, \textit{EPYC}, \textit{LUM}, and \textit{OGN}. These genes were not selected
because RNA abundance alone demonstrates function. They represent several SLRP
classes and combine juvenile-lineage expression with different degrees of
growth-plate localization, skeletal function, and collagen-matrix evidence.
The literature also makes clear that SLRPs are partly redundant: a mild
single-gene phenotype can coexist with a stronger compound-knockout phenotype,
whereas a compound knockout cannot identify the separate contribution of each
protein \citep{Nikitovic2012}.

\subsubsection{Biglycan, decorin, and fibromodulin}

\textit{BGN} encodes biglycan, a class-I chondroitin/dermatan-sulphate SLRP
abundant in skeletal connective tissues. Biglycan contributes to collagen-fibril
organization, osteogenesis, and bone remodelling. Biglycan-deficient mice show
abnormal collagen fibrils and reduced skeletal growth or bone mass, with a much
stronger phenotype after combined loss of the related class-I SLRP decorin
\citep{Corsi2002}. Biglycan therefore has both a structural matrix role and
context-dependent effects on bone-cell signalling; the relative contribution of
these mechanisms within individual growth-plate zones remains unresolved.

\textit{DCN} encodes decorin, a closely related class-I SLRP with established
roles in collagen fibrillogenesis, cartilage matrix organization, and cell-
matrix signalling. It was included both as a biologically relevant candidate
and as the essential paralog control for testing whether apparent BGN sequences
were actually decorin-like records.

\textit{FMOD} encodes fibromodulin, a class-II collagen-binding SLRP. Loss of
fibromodulin disrupts tendon collagen-fibril morphology and alters lumican
deposition, demonstrating an in-vivo role in collagen fibrillogenesis and
partial functional overlap with other SLRPs \citep{Svensson1999}. During mouse
knee development, fibromodulin is enriched in growing epiphyses and localizes
around late-hypertrophic chondrocytes, the secondary ossification centre, and
the growth plate \citep{Saamanen2001}. FMOD is also cleaved by cartilage-
remodelling proteases, and MMP-13-cleaved fibromodulin occurs around
hypertrophic chondrocytes in fetal growth plates \citep{Shu2019}.

BGN and FMOD also influence skeletal remodelling together. Their combined loss
in mice increases osteoclastogenesis and produces low bone mass, while both
proteins can interact with TNF-alpha and RANKL in experimental systems
\citep{Kram2017}. However, incomplete rescue of the double-knockout phenotype by
OPG-Fc, together with adverse effects in young mice, shows that this biology
cannot be reduced to a single osteoclast pathway \citep{Kram2020}.

\subsubsection{PRELP}

\textit{PRELP} encodes a class-II extracellular LRR protein distinguished by a
basic proline/arginine-rich N-terminus. This N-terminal region binds heparin and
heparan sulphate and may link matrix components to cell-surface heparan-sulphate
proteoglycans \citep{Bengtsson2000}. PRELP is present in cartilage, and its LRR,
cysteine, and glycosylation-associated regions are conserved between human and
mouse \citep{Grover1996,Grover2002}. PRELP-derived N-terminal peptides inhibit
osteoclast differentiation and can reduce bone loss in experimental models,
whereas cell-culture work also supports a role in osteoblast differentiation
\citep{Rucci2009,Rucci2013,Li2016PRELP}. Peptide pharmacology and cultured-cell
effects do not yet establish the normal function of endogenous full-length PRELP
in the growth plate, so its precise in-vivo mechanism remains incomplete.

\subsubsection{Epiphycan}

\textit{EPYC} encodes epiphycan/PG-Lb, a class-III cartilage-enriched
chondroitin/dermatan-sulphate SLRP. In embryonic mouse limbs, \textit{Epyc} RNA
is associated mainly with resting and proliferating chondrocytes, while the
secreted protein is distributed throughout the growth-plate matrix, including
around hypertrophic chondrocytes \citep{Johnson1999}. Epyc-null mice are viable
and lack a severe early skeletal phenotype, but develop age-related joint and
femur-length abnormalities; earlier or stronger osteoarthritis in
Epyc/Bgn-deficient mice indicates interaction and compensation between SLRPs
\citep{Nuka2010}. EPYC therefore has direct anatomical relevance to the growth
plate, but its specific molecular binding partners and indispensable mechanism
remain less well defined than its distribution.

\subsubsection{Lumican}

\textit{LUM} encodes lumican, a widely distributed class-II collagen-binding
SLRP. Lumican regulates collagen-fibril diameter and organization in several
connective tissues and has partly overlapping functions with fibromodulin.
During skeletal development it is present in cartilaginous and developing bone
matrices \citep{Raouf2002}. Lumican knockdown in tissue-engineered cartilage
increased type-II collagen deposition and fibril diameter, demonstrating that
its effect on collagen accumulation is context dependent rather than simply
``pro-collagen'' \citep{Kafienah2008}. Recombinant lumican also suppresses
osteoclast differentiation and resorption through reduced Akt activity in
vitro \citep{Lee2021Lum}; whether this mechanism operates at physiologically
important levels in the juvenile growth plate remains unknown.

\subsubsection{Osteoglycin}

\textit{OGN}, also known as mimecan, encodes a class-III secreted SLRP. Ogn-null
mice are viable without an obvious gross developmental phenotype, but show
abnormal collagen fibril diameter and morphology in skin and cornea, providing
in-vivo evidence for collagen-fibrillogenesis control \citep{Tasheva2002}.
BMP1/Tolloid-family metalloproteinases process pro-OGN and enhance its regulation
of type-I collagen fibrillogenesis in biochemical assays \citep{Ge2004}.
Vertebrate OGN history also includes teleost duplication and
subfunctionalization \citep{Costa2018}. OGN therefore contributes an informative
evolutionary and matrix-remodelling case, but it has weaker direct evidence for
a growth-plate-specific function than FMOD or EPYC.

\subsubsection{Osteomodulin as a context gene}

\textit{OMD}/osteoadherin is not a seventh fully analysed panel member, but it
provides relevant cartilage--bone-interface context. The secreted protein binds
hydroxyapatite, supports integrin-associated osteoblast attachment, and is
localized to the primary spongiosa of fetal growth plates \citep{Sommarin1998}.
OMD has also been linked to BMP2/SMAD-dependent osteogenesis
\citep{Lin2021OMD}. Its expression and representative gene structure are
reported here, but it did not undergo the complete protein, phylogenetic, and
synteny workflow.

Together, these studies justify the panel as a set of biologically plausible
cartilage and skeletal ECM genes while also defining the limits of the thesis.
The comparative analyses can test ortholog plausibility and conservation of
secretion, LRR architecture, sequence, gene structure, phylogenetic placement,
and genomic neighbourhood. They cannot demonstrate that every ortholog has the
same collagen-binding, signalling, or growth-plate phenotype.

\subsection{Expression-guided evolutionary analysis}
\label{subsec:expression-guided-evolution}

Because SLRPs form a multi-gene family with related domain architectures and partially overlapping functions, sequence similarity alone may not be sufficient to distinguish orthologs from closely related paralogs. Combining expression evidence with reciprocal sequence searches, domain validation, phylogenetic reconstruction, and synteny analysis can therefore provide a more reliable framework for identifying conserved cartilage-associated SLRP candidates across vertebrates.

\subsection{Aim of this thesis}
\label{subsec:aim}

This thesis aims to characterize selected growth-plate-associated SLRP genes by
combining juvenile and anatomically resolved expression evidence with
comparative analyses across representative vertebrates. Specifically, it asks
which SLRPs form a biologically justified candidate panel, whether their
secretory signals, LRR-rich protein cores, coding organization, phylogenetic
identity, and genomic neighbourhoods are conserved, and where deep-lineage or
annotation uncertainty limits an ortholog assignment. By integrating these
independent evidence layers, the work distinguishes a conserved vertebrate ECM
framework from candidate-specific expression differences, paralogy, and gene-
model uncertainty.
